\section{Unidad territorial y universo de análisis}
\label{ch:unidad_analisis}

\section{Municipio como unidad de análisis}
\label{sec:municipio_unidad}

El municipio es la unidad territorial básica de la organización político-administrativa colombiana (Constitución de 1991, Art. 286; Ley 136 de 1994). Es el nivel donde convergen:

\begin{itemize}
    \item Competencias en prestación de servicios públicos (educación, salud, agua, saneamiento).
    \item Ejecución de contratación pública local y departamental.
    \item Disponibilidad de información censal y estadística oficial desagregada.
    \item Identidad territorial para análisis de vulnerabilidad y políticas públicas.
\end{itemize}

La elección del municipio (y no el departamento, la localidad o el corregimiento) responde a la necesidad de \textbf{comparabilidad nacional} y \textbf{coherencia con la unidad de contratación pública} (SECOP II registra entidad contratante y lugar de ejecución a nivel municipal).

\section{Códigos territoriales}
\label{sec:codigos_territoriales}

\begin{itemize}
    \item \textbf{Formato:} \texttt{codigo\_municipio} como \textbf{texto de 5 caracteres} (DIVIPOLA), p.ej. ``11001'' (Bogotá), ``05001'' (Medellín).
    \item \textbf{Nunca se usa el nombre municipal como llave} (riesgo de homónimos, variaciones ortográficas, cambios de nombre).
    \item \textbf{Longitud fija:} Se valida que todos los códigos tengan 5 caracteres; se completan con ceros a la izquierda si es necesario.
    \item \textbf{Consistencia entre fuentes:} La unión entre tablas CNPV, IPM y DIVIPOLA se realiza exclusivamente por \texttt{codigo\_municipio}.
\end{itemize}

\section{Cobertura municipal}
\label{sec:cobertura_municipal}

El universo de trabajo adoptado son los \textbf{1.122 códigos municipales} que tienen datos en el CNPV 2018 y el IPM 2018. Esta cifra se valida en \texttt{resultados/base\_trabajo/validaciones.csv} (validación ``Llave/prefijo ipm\_municipal\_ancho'': 1.122 evaluados, 0 violaciones).

\begin{table}[htbp]
\centering
\caption{Universo municipal y cobertura}
\label{tab:universo_municipal}
\begin{tabular}{@{}lrr@{}}
\toprule
\textbf{Concepto} & \textbf{Cantidad} & \textbf{Fuente/Evidencia} \\
\midrule
Municipios en CNPV 2018 (cuadros municipales) & 1.122 & \texttt{universo\_municipal\_maestro.csv} \\
Municipios en IPM 2018 (hoja municipal) & 1.122 & \texttt{validaciones.csv} \\
Municipios en DIVIPOLA 20260917 & $\approx$1.122 & \texttt{DIVIPOLA-\_Codigos\_municipios\_20260917.csv} \\
Coincidencia CNPV $\cap$ IPM & 1.122 & \texttt{Final\_Base.md}, validaciones \\
\bottomrule
\end{tabular}
\end{table}

\section{Coincidencias entre fuentes}
\label{sec:coincidencias}

La homologación territorial entre CNPV 2018, IPM 2018 y DIVIPOLA se realizó mediante \texttt{codigo\_municipio}. Los controles de llave/prefijo en \texttt{validaciones.csv} reportan \textbf{cero violaciones} para todas las tablas intermedias (\texttt{contexto\_1PM}, \texttt{contexto\_1V}, \texttt{hogares\_1HM\_2HM}, \texttt{persona\_demo\_17PM}, \texttt{persona\_social\_5PM}, \texttt{viviendas\_6V}, \texttt{viviendas\_8V}, \texttt{viviendas\_10V}, \texttt{ipm\_municipal\_ancho}).

\section{No coincidencias y casos especiales}
\label{sec:no_coincidencias}

A pesar de la coincidencia en el conteo total (1.122), existen \textbf{diferencias estructurales} que no son errores sino consecuencias de definiciones territoriales distintas:

\subsection{Municipios sin cabecera (20 casos)}
\label{sec:municipios_sin_cabecera}

Veinte municipios no tienen área ``Cabecera'' censada en 2018. Estos coinciden exactamente con los \textbf{20 valores \texttt{NA} estructurales del IPM para Cabeceras} (validación ``20 NA IPM cabeceras = municipios sin cabecera'': 20 evaluados, 0 violaciones).

\begin{itemize}
    \item 18 corresponden a \textbf{Áreas No Municipalizadas (ANM)} en DIVIPOLA 2026.
    \item 2 códigos adicionales: \texttt{94343} (Barrancominas, 7.207 hab.) y \texttt{94663} (Mapiripana, 1.011 hab.), cuya población rural dispersa suma 8.218 hab.
    \item Población total en los 20 municipios sin cabecera: \textbf{34.787 habitantes} (incluye 26.569 en las 18 ANM + 8.218 en los dos códigos adicionales).
    \item Proporción rural dispersa en estos municipios: 47,27\% -- 100\%.
    \item Detalle en \texttt{resultados/base\_trabajo/municipios\_sin\_cabecera.csv} y \texttt{resumen\_sin\_cabecera\_departamento.csv}.
\end{itemize}

\subsection{Código 94663 (Mapiripana) -- Relación con 94343 (Barrancominas)}
\label{sec:94663}

\begin{itemize}
    \item Documentado para la capa SECOP futura: \texttt{94663 $\rightarrow$ 94343}.
    \item En los datos CNPV/IPM 2018 \textbf{permanecen ambos códigos separados}; no se fusionan ni reasignan conteos.
    \item Ver \texttt{Final\_Base.md} y \texttt{crosswalk\_enlace.csv}.
\end{itemize}

\subsection{Código 27493 (Nuevo Belén de Bajirá) -- Conflicto territorial}
\label{sec:27493}

\begin{itemize}
    \item Se propone \texttt{27615} (Riosucio) como ``padre'' territorial, \textbf{pendiente de verificación DANE/IGAC}.
    \item Existe conflicto histórico con \texttt{05480} (Mutatá).
    \item Evidencia compatible: Riosucio suma 48.257 habitantes y un centro poblado de 13.504; Mutatá suma 12.607. No prueba el límite; \textbf{no se fraccionan ni reasignan conteos CNPV}.
    \item Ver \texttt{Final\_Base.md} y \texttt{crosswalk\_enlace.csv}.
\end{itemize}

\subsection{Municipios sin centro poblado (192 casos)}
\label{sec:sin_centro_poblado}

192 municipios no tienen área ``Centro poblado'' en 2018. Se conservan como \textbf{cobertura territorial estructural}; no son faltantes estadísticos. Ver \texttt{Final\_Base.md} (validación ``Municipios sin Centro poblado'': 192, PASS).

\subsection{Áreas territorialmente inexistentes}
\label{sec:areas_inexistentes}

El archivo \texttt{resultados/base\_trabajo/areas\_inexistentes.csv} documenta, por municipio, qué áreas (Cabecera, Centro poblado, Rural disperso) no existen. Esto permite distinguir \textbf{ausencia estructural} (el área no existe en el territorio) de \textbf{faltante estadístico} (el dato no se capturó). Ejemplo: en un municipio 100\% urbano, la ausencia de ``Rural disperso'' es estructural, no un dato que deba imputarse.

\subsection{Ausencias estructurales en combinaciones demográficas}
\label{sec:ausencias_estructurales}

\begin{itemize}
    \item \textbf{Grupos de edad:} 335 combinaciones municipio-área-grupo de edad ausentes en 1PM (119 municipios afectados). Ver \texttt{ausencias\_estructurales\_1PM.csv}.
    \item \textbf{Tipo de vivienda:} 1.202 combinaciones ausentes.
    \item \textbf{Tratamiento:} Se completan como \textbf{ceros estructurales} y se marcan con \textit{flags}; \textbf{no se imputan} estadísticamente. No se crean filas para áreas inexistentes.
\end{itemize}

\section{Resumen de decisiones territoriales}
\label{sec:resumen_territorial}

\begin{table}[htbp]
\centering
\caption{Decisiones territoriales clave}
\label{tab:decisiones_territoriales}
\begin{adjustbox}{max width=\textwidth}
\begin{tabular}{@{}lll@{}}
\toprule
\textbf{Aspecto} & \textbf{Decisión} & \textbf{Evidencia} \\
\midrule
Universo municipal & 1.122 códigos CNPV/IPM 2018, sin fusiones & \texttt{Final\_Base.md}, \texttt{universo\_municipal\_maestro.csv} \\
Llave de unión & \texttt{codigo\_municipio} (texto 5 chars) & \texttt{R/01\_verificar\_consolidados.R} \\
Municipios sin cabecera (20) & Conservados; NA estructural en IPM Cabecera & \texttt{municipios\_sin\_cabecera.csv} \\
Código 94663 & Relación futura 94663 $\rightarrow$ 94343 (solo enlace SECOP) & \texttt{crosswalk\_enlace.csv} \\
Código 27493 & Hipótesis 27615 como padre; conflicto con 05480; sin reasignar & \texttt{crosswalk\_enlace.csv} \\
Ausencias estructurales & Ceros estructurales con flags; no imputación & \texttt{ausencias\_estructurales\_1PM.csv}, \texttt{areas\_inexistentes.csv} \\
Cobertura 5PM / pob\_10mas & 28 municipios con cobertura $<0,9$ (umbral configurable) & \texttt{validaciones.csv}, \texttt{R/00\_config.R} \\
\bottomrule
\end{tabular}
\end{adjustbox}
\end{table}
